\section{Advanced worked circuits}
\hypertarget{advanced-circuits}{}

The earlier examples use one or two gates. The circuits in this chapter show
that the same pieces scale to arithmetic, comparison, parity, phase tricks,
and small oracle algorithms. Every listing is a \textbf{Complete runnable
protocol}: paste it into \textbf{Compile text into a circuit}, choose
\textbf{Compile protocol}, then \textbf{Run all}. Circuits with \code{PARAMS}
can also be checked row by row in the \textbf{Circuit correction lab} with
\textbf{Infer truth table}, \textbf{Probe outputs}, and \textbf{Test circuit}.

\begin{beginnerbox}[One pattern runs through this whole chapter]
Inputs are \emph{read}, never overwritten. Each answer is written into a
fresh output wire that starts at 0, using the rule
$\text{target}' = \text{target}\oplus f(\text{inputs})$. When a circuit needs a
temporary helper wire, it computes the helper, uses it, and then runs the same
gates again to return the helper to 0 (\emph{uncomputation}).
\end{beginnerbox}

\subsection{Full adder from primitive gates}

A full adder adds three bits $A$, $B$, and a carry-in $C_{\rm in}$:
\[
\mathrm{Sum}=A\oplus B\oplus C_{\rm in},\qquad
C_{\rm out}=\operatorname{majority}(A,B,C_{\rm in})
=(A\land B)\oplus(A\land C_{\rm in})\oplus(B\land C_{\rm in}).
\]

\textbf{Complete runnable protocol}
\begin{lstlisting}
PARAMS: A:state B:state Cin:state

MAIN-PROCESS FlatFullAdder
CREATETOKEN -I Cout Sum
SET Cout 0p
SET Sum 0p

# Stage 1: Sum collects the parity of the three inputs.
CNOT -I A -O Sum
CNOT -I B -O Sum
CNOT -I Cin -O Sum

# Stage 2: Cout flips once for every pair of inputs that are both 1.
CCNOT -I A B -O Cout
CCNOT -I A Cin -O Cout
CCNOT -I B Cin -O Cout

MEASURE -I Cout
MEASURE -I Sum
RETURNVALS Cout Sum
\end{lstlisting}

\textbf{Why the carry works:} if exactly two inputs are 1, one pair is active
and Cout flips once, to 1. If all three are 1, all three pairs are active and
Cout flips three times, which still leaves 1. With zero or one active input no
pair fires. That is exactly ``at least two of three'', the majority function.
The bundled \code{SingleBitFullAdder} is the same circuit.

\subsection{Two-bit adder built from a child process}

Multi-bit addition chains full adders: the carry out of bit 0 becomes the
carry in of bit 1. \code{RUNCHILD} reuses the bundled \code{SingleBitFullAdder}
for each bit, so the parent only describes the wiring.

\textbf{Complete runnable protocol}
\begin{lstlisting}
PARAMS: A1:state A0:state B1:state B0:state

MAIN-PROCESS TwoBitRippleAdder
DECLARECHILD SingleBitFullAdder
CREATETOKEN -I NoCarry C0 S0 C1 S1
SET NoCarry 0p

# Bit 0: A0 + B0 with no incoming carry.
RUNCHILD SingleBitFullAdder -I A0 B0 NoCarry -O C0 S0
# Bit 1: A1 + B1 plus the carry from bit 0.
RUNCHILD SingleBitFullAdder -I A1 B1 C0 -O C1 S1

MEASURE -I C1
MEASURE -I S1
MEASURE -I S0
RETURNVALS C1 S1 S0
\end{lstlisting}

The outputs read as a three-bit binary number $C_1S_1S_0$. For example,
$A=11_2=3$ and $B=10_2=2$ give $C_1S_1S_0=101_2=5$. The bundled
\code{TwoBitFullAdder} and \code{FourBitFullAdder} extend the same chain with
an explicit carry input and more bits.

\subsection{Reversible comparator: is A greater than B?}

For single bits, $A>B$ exactly when $A=1$ and $B=0$, that is
$A\land\neg B$. The circuit negates B temporarily, uses it, and restores it.

\textbf{Complete runnable protocol}
\begin{lstlisting}
PARAMS: A:state B:state

MAIN-PROCESS GreaterThan
CREATETOKEN -I Gt
SET Gt 0p

X -I B -O B
AND -I A B -O Gt
X -I B -O B

MEASURE -I Gt
RETURNVALS Gt
\end{lstlisting}
\begin{center}
\begin{tabular}{cc|c}\toprule $A$&$B$&Gt\\\midrule
0&0&0\\0&1&0\\1&0&1\\1&1&0\\\bottomrule\end{tabular}
\end{center}
The second X is what keeps the circuit honest: without it, B would leave the
circuit inverted and any later gate reading B would see the wrong value.

\subsection{Two-bit equality with uncomputation}

Two numbers are equal when every bit agrees. This circuit writes a
``this bit differs'' flag for each bit, combines the flags, and then cleans the
flags up.

\textbf{Complete runnable protocol}
\begin{lstlisting}
PARAMS: A1:state A0:state B1:state B0:state

MAIN-PROCESS TwoBitEquality
CREATETOKEN -I D1 D0 Eq
SET D1 0p
SET D0 0p
SET Eq 1p

# Compute: D = 1 where the bits differ.
XOR -I A1 B1 -O D1
XOR -I A0 B0 -O D0

# Eq starts at 1 and flips if any bit differs.
OR -I D1 D0 -O Eq

# Uncompute: the same XORs return D1 and D0 to 0.
XOR -I A1 B1 -O D1
XOR -I A0 B0 -O D0

MEASURE -I Eq
MEASURE -I D1
MEASURE -I D0
RETURNVALS Eq D1 D0
\end{lstlisting}

Three ideas meet here:
\begin{description}
  \item[A target that starts at 1] \code{SET Eq 1p} makes
  $\mathrm{Eq}'=1\oplus(D_1\lor D_0)=\neg(D_1\lor D_0)$. Starting a target
  at 1 turns ``flip when true'' into ``clear when true''.
  \item[Uncomputation] XOR is its own inverse, so repeating the two XOR lines
  returns $D_1$ and $D_0$ to 0 for every input.
  \item[Proof by truth table] Returning D1 and D0 lets a truth-table test
  confirm the cleanup. In all 16 rows those two columns must be 0.
\end{description}

\subsection{Parity of four bits}

Parity is 1 when an odd number of inputs are 1. Each CNOT adds one input into
the target modulo 2.

\textbf{Complete runnable protocol}
\begin{lstlisting}
PARAMS: A:state B:state C:state D:state

MAIN-PROCESS FourBitParity
CREATETOKEN -I P
SET P 0p

CNOT -I A -O P
CNOT -I B -O P
CNOT -I C -O P
CNOT -I D -O P

MEASURE -I P
RETURNVALS P
\end{lstlisting}
Parity checks are how error-detecting codes notice a single flipped bit. The
two-input \code{XOR -I A B -O P} is shorthand for the first two CNOT lines.

\subsection{Phase kickback}
\hypertarget{phase-kickback}{}

A CNOT never changes its control's 0/1 value, but it \emph{can} change the
control's phase. Put the target in $\ket-=(\ket0-\ket1)/\sqrt2$: flipping
$\ket-$ gives $-\ket-$, so ``flip the target'' becomes ``multiply by $-1$'',
and that sign attaches to the control's $\ket1$ branch.

\textbf{Complete runnable protocol}
\begin{lstlisting}
MAIN-PROCESS PhaseKickback
CREATETOKEN -I Ctrl Tgt
SET Ctrl 0p
SET Tgt 1p

H -I Ctrl -O Ctrl
H -I Tgt -O Tgt
CNOT -I Ctrl -O Tgt
H -I Ctrl -O Ctrl

MEASURE -I Ctrl
RETURNVALS Ctrl
\end{lstlisting}
\[
\ket{+}\ket{-}\xrightarrow{\mathrm{CNOT}}\ket{-}\ket{-}
\xrightarrow{H\text{ on Ctrl}}\ket{1}\ket{-}.
\]
Ctrl measures 1 every time. Change \code{SET Tgt 1p} to \code{SET Tgt 0p}
and it measures 0 every time: the target's state decided the control's
result, even though the control was never ``written''. The \emph{Phase
kickback} starter circuit shows the same effect on the canvas.

\subsection{A small oracle: finding a hidden bit string}

An \emph{oracle} is a circuit treated as a black box that computes some
$f(x)$. Here $f(x)=s\cdot x \bmod 2$ for a hidden three-bit string $s$
(Bernstein--Vazirani). Classically, finding $s$ takes three queries. Using
kickback, one query is enough.

\textbf{Complete runnable protocol}
\begin{lstlisting}
MAIN-PROCESS HiddenString101
CREATETOKEN -I In2 In1 In0 Anc
SET In2 0p
SET In1 0p
SET In0 0p
SET Anc 1p

H -I In2 -O In2
H -I In1 -O In1
H -I In0 -O In0
H -I Anc -O Anc

# Oracle for s = 101: one CNOT for every 1 in s.
CNOT -I In2 -O Anc
CNOT -I In0 -O Anc

H -I In2 -O In2
H -I In1 -O In1
H -I In0 -O In0

MEASURE -I In2
MEASURE -I In1
MEASURE -I In0
RETURNVALS In2 In1 In0
\end{lstlisting}
The measured outputs are $1,0,1$: the hidden string itself. Each oracle
CNOT kicks a phase back onto one input wire, and the final H gates turn
those phases into 0/1 values. Change which inputs have an oracle CNOT and the
answer changes to match. With a single input wire this is Deutsch's
algorithm.

\subsection{Two-qubit Grover search}

Grover's algorithm finds a marked input by boosting its amplitude. With two
qubits, one round finds the answer with certainty.

\textbf{Complete runnable protocol}
\begin{lstlisting}
MAIN-PROCESS GroverFind11
CREATETOKEN -I A B
SET A 0p
SET B 0p

# Equal superposition of 00, 01, 10, 11.
H -I A -O A
H -I B -O B

# Oracle: mark 11 with a minus sign.
CZ -I A -O B

# Diffusion: reflect every amplitude about the average.
H -I A -O A
H -I B -O B
X -I A -O A
X -I B -O B
CZ -I A -O B
X -I A -O A
X -I B -O B
H -I A -O A
H -I B -O B

MEASURE -I A
MEASURE -I B
RETURNVALS A B
\end{lstlisting}
\begin{center}
\begin{tabular}{l|cccc}\toprule
Amplitude of & $\ket{00}$ & $\ket{01}$ & $\ket{10}$ & $\ket{11}$\\\midrule
After H, H & $\tfrac12$ & $\tfrac12$ & $\tfrac12$ & $\tfrac12$\\
After oracle & $\tfrac12$ & $\tfrac12$ & $\tfrac12$ & $-\tfrac12$\\
After diffusion & 0 & 0 & 0 & $\pm1$\\\bottomrule
\end{tabular}
\end{center}
The oracle does not change any probability, since every entry still has
magnitude $\tfrac12$. The diffusion stage turns the marked sign into certainty.
Open \textbf{Outcome truth table / probabilities} on the
\textbf{Particle visualization} page after each \textbf{Step gate} to watch
the amplitudes move.
